% ============================================================================
% DARPA DSO Office-wide BAA HR001126S0015: Executive Summary
% (Attachment A template, mandatory for executive summaries)
%
% Format (Attachment A, Executive Summary Template, pp. 1-3):
%   - PDF or Word; English; 8.5 x 11 in paper; 1 in margins
%   - font no smaller than 12 pt (8 or 10 pt allowed in figures, tables, charts)
%   - at most 2 pages: 1 page cover sheet + 1 page technical description
%   - cover sheet: prime organization logo, executive summary title, DSO thrust
%     area, proposer organization, technical POC and administrative POC (name,
%     mailing address, telephone, email)
%   - technical description: answers to four Heilmeier questions
% Labeling (DARPA Proposer Instructions: General Terms and Conditions, "Content
% and Form of Application Submission"): every document carries the BAA number,
% proposer organization,
% and title or short title. The running header does that on both pages.
% Thrust areas named as in the BAA, Section I.A, p. 3.
%
% Evidence numbers come only from afosr-yip-white-paper-a2c.tex, Section 5,
% and benchmarks/bofn-prelim/README.md.
%
% Build:  xelatex executive-summary.tex  (one pass is enough)
% Lint:   cringe-filter lint -c proposal executive-summary.tex
% ============================================================================

\documentclass[12pt,letterpaper]{article}

\usepackage[margin=1in,headheight=15pt,headsep=0.3in]{geometry}

\usepackage{iftex}
\ifPDFTeX
  \usepackage[T1]{fontenc}
  \usepackage[utf8]{inputenc}
  \usepackage{mathptmx}
\else
  \usepackage{fontspec}
  \IfFontExistsTF{Times New Roman}
    {\setmainfont{Times New Roman}}
    {\setmainfont{TeX Gyre Termes}}
\fi
\usepackage{microtype}
\usepackage{xcolor}
\usepackage{array}
\usepackage[colorlinks=true,urlcolor=blue!70!black,linkcolor=blue!70!black]{hyperref}

\setlength{\parindent}{0pt}
\setlength{\parskip}{4pt plus 1pt minus 1pt}

% Running header: BAA number, organization, short title (12 pt).
\usepackage{fancyhdr}
\pagestyle{fancy}
\fancyhf{}
\renewcommand{\headrulewidth}{0pt}
\fancyhead[L]{HR001126S0015}
\fancyhead[C]{Brigham Young University}
\fancyhead[R]{Choosing the Next Experiment}
\fancyfoot[C]{\thepage}

% Red notes for Sterling; remove all before submission.
\newcommand{\todo}[1]{\textcolor{red}{[TODO Sterling: #1]}}

\newcommand{\q}[1]{\textbf{#1}\enspace}

\begin{document}

% ===========================================================================
% Page 1: cover sheet
% ===========================================================================
\begin{center}
\fbox{\parbox[c][0.9in][c]{2.8in}{\centering
  \todo{insert the official BYU logo here}}}

\vspace{14pt}
{\bfseries COVER SHEET\par}
\vspace{4pt}
Executive Summary, DARPA Defense Sciences Office (DSO)\\
Office-wide Broad Agency Announcement HR001126S0015
\end{center}

\vspace{10pt}
\renewcommand{\arraystretch}{1.25}
\begin{tabular}{@{}>{\bfseries\raggedright}p{2.0in}>{\raggedright\arraybackslash}p{4.3in}@{}}
Executive Summary Title &
  Choosing the Next Experiment: Finding Printable Alloys at One Tenth
  the Cost \\[4pt]
DSO Thrust Area &
  (1) Materials, Manufacturing, and Energy Systems;\newline
  (4) Complex, Dynamic, and Intelligent Systems \\[4pt]
Proposer Organization &
  Brigham Young University (BYU), Provo, Utah \\[10pt]
\multicolumn{2}{@{}l}{\bfseries Technical Point of Contact (POC)} \\
\quad Name & Sterling G.\ Baird, Ph.D., Assistant Professor,
  Department of Mechanical Engineering \\
\quad Mailing Address & Brigham Young University, 350 EB, Provo, UT 84602 \\
\quad Telephone & (801) 422-5344 \\
\quad Email & \href{mailto:sterling.baird@byu.edu}{sterling.baird@byu.edu}
  \\[10pt]
\multicolumn{2}{@{}l}{\bfseries Administrative POC} \\
\quad Name & \todo{name and title of the BYU research administration
  contact} \\
\quad Mailing Address & \todo{mailing address} \\
\quad Telephone & \todo{telephone} \\
\quad Email & \todo{email} \\
\end{tabular}

\clearpage

% ===========================================================================
% Page 2: technical description (one page)
% ===========================================================================
{\bfseries Executive Summary Technical Description\par}
\vspace{2pt}

\q{What is the proposed work attempting to accomplish or do?}
We aim to cut the cost and time of finding and qualifying a new printable
alloy by about ten times, and to do the same for other multi-step materials
workflows. Developing an alloy for laser powder bed fusion (LPBF), a form of
metal 3D printing, takes five steps: a thermodynamic
calculation (minutes), powder production (a day), printed tracks and cubes
(hours), electron microscopy (a day), and tensile tests (days). We will
build and prove a method that decides, as each result arrives, which step to
run next and on which candidate alloy.

\q{How is the work performed today, and what are the limitations?}
Experimental campaigns usually send every candidate through every step, so
much of the budget goes to full tests of alloys that a cheap early
measurement could have ruled out. Computational campaigns work as a funnel:
many simulations, then a few survivors go to experiments. The funnel is only
as good as its physics. Standard thermodynamic screens miss effects such as
zirconium additions that keep printed high-strength
aluminum from cracking (Martin et al.,
\emph{Nature} 2017), so the funnel can discard the best alloys. Published
methods that learn from partial results optimize a single property, wait
for each result before choosing the next experiment, and have been tested only on synthetic problems.

\q{What is new in our approach, and why will it succeed?}
Each step gets its own statistical model, and the models are linked in the
order of the steps, so uncertainty carries through the chain (Astudillo and
Frazier, NeurIPS 2021). For each possible experiment, the algorithm divides
the expected gain by the cost and runs the best one, so a candidate can stop
at any step (Buathong et al., ICML 2024). Simulations run continuously at full
throughput on BYU's computing cluster, and each choice uses all data available at that
moment. We will add three things not done before: balancing several
properties at once (freedom from cracks, strength, ductility); choosing new
experiments while others are still running; and a demonstration on a real
five-step workflow with build-to-build scatter. In a hand-built simulation of
this chain whose crack screen ignores zirconium, the screening funnel
stalled at 63\% of the best possible alloy after \$100,000. The algorithm
reached that level after about \$26,500, or 3.8 times less spend (90\%
confidence interval 2.0 to 5.7; 20 campaigns). When the screen's physics was
right, screening kept pace. The simulation is noise-free, with one objective
and one experiment at a time, so ten times is our target, not yet a
result. The
full chain already runs for aluminum alloys: our ultrasonic atomizer makes
powder of four compositions per week, LPBF builds cost \$800 each at a
University of Utah facility, and the BYU Microscopy Center supplies electron
microscopy. We have completed a powder-to-part demonstration on a baseline
aluminum alloy. The PI developed Honegumi, open software that helps
experimentalists set up such searches, and co-authored the
\emph{Chemical Reviews} survey of self-driving laboratories.

\q{Who will care, and what will the impact be if the work is successful?}
Qualifying a new material for a Department of War (DoW) system takes many
expensive tests in a fixed order. A method that picks the next test and
states its confidence shortens that path for any multi-step materials or
manufacturing workflow. The first target is high-strength aluminum alloys
that print without cracking, for aircraft structures and for replacement
parts no longer made. The method moves to a new workflow by swapping in its
steps and their costs. A DARPA-scale effort would add a plasma melting module
for our atomizer (about \$150,000), extending the chain to refractory alloys
for hypersonic structures, and more printer capacity to run many builds at
once. We will release software, data, and models openly for DoW
laboratories to use on their own materials.

\end{document}
